
\chapterimage{livingmtl.pdf}
\chapterimagecaption{\textit{A Sunday Afternoon on the Island of La Grande Jatte} (1886) by Georges Seurat}

\chapter{Living in Montreal}
\section{Housing and Transportation}
Here we have compiled some resources that you may find useful while living in Montreal as a graduate student including information on accommodation, local transit, bike and car rental.

\subsection{Accommodation}\index{Accommodation}
Montréal has significant \href{https://www.centris.ca/en/blog/real-estate/average-price-of-apartments-in-montreal-in-2025}{price differences between neighbourhoods}. As you might expect, the further from the downtown and the less accessible by public transit a sector is, the cheaper the rent. With the addition of the REM, transit times to McGill have been significantly reduced and new places are now viable for a lower price! You should look into the neighbourhood and the apartment alike. Some parts of the city make it a much nicer experience, while others can also make it worse. 

\noindent While finding apartments through official real estate agencies is common, it is typical to use \textbf{Facebook Marketplace}. Usually, Marketplace has landlords with fewer apartments. Some people prefer the human connection of knowing the landlord, while others prefer the administrative structure and efficiency of companies. In either case, it is a good idea to look into each neighbourhood (see previous point) and learn a bit about your rights and local laws (see following point). Lease transfers are also a good way to find rentals at a good quality to price ratio.

\begin{plainexercise}
Landlords are only allowed to ask for an upfront payment of the first month of rent. Deposits are illegal. 
\end{plainexercise} \vspace{-5pt}

\noindent While not a complete list by any means, here are some more important details for rent/leases: 
\begin{plaincorollary}
    \begin{itemize}[leftmargin=10pt]
        \item While your landlord is required to send you notification should any changes be made to your lease prior to renewal, the lease itself \textbf{auto-renews} if no action is taken. If your landlord sends you a renewal notice, you have a month to express that you will not renew the lease. Otherwise, you automatically renew with the modified conditions. If you do not receive a notice, you still have to express your intent to move out within 3-6 months prior to the end of your lease lest the current lease is extended for an additional 12 months. 
        \item  Rent increases in Québec are well regulated by the Tribunal administratif du logement (TAL). There is no absolute legal rent cap, but increases must be justified by building operating costs, taxes, and inflation. Every year, the TAL publishes a \href{https://www.tal.gouv.qc.ca/en/calculation-for-rent-increase}{calculation tool} for allowable rent increase as well as the recommended increase percentage. 
        \item While not law-related, most apartments come unfurnished and without appliances (washer, dryer, fridge, oven, etc.). You should check with the landlord or current tenants to see if such things are included. 
        \item Here is a website that summarises the main points of Québec law: \href{https://educaloi.qc.ca/en/web-guide/rental-housing/}{Éducaloi}. If you have doubts, it might be better to look it up online.
    \end{itemize} 
\end{plaincorollary}

\subsection{Transit} \index{Transit}
As a graduate student at McGill you are eligible for the Student Opus card which allows you to load student fares onto your card that are almost 50\% cheaper than regular adult fares. The student card costs \$15 and needs to be ordered online via Minerva.

\textit{Minerva > Student Menu > Reduced-fare OPUS card order form}.

The \href{https://www.stm.info/sites/default/files/pdf/en/tarifs.pdf}{monthly fare} for students in \href{https://www.artm.quebec/zones-tarifaires/}{Zone A} is \$66. Instead, buying a 4-month plan for \$251 will save you \$13. 

 If you ever go to an evening event outside of Zone A or if you don't have a transit pass, the STM evening pass might save you some money. It only works between 6 pm and 5 am, but it is cheaper than 2 tickets, works for all zones, and has unlimited uses during that time!

\subsection{Biking} \index{Biking}
If you don't plan to use the metro or buses much during the summer or if you want to take advantage of the (usually) nice weather to get some exercise in, the BIXI service in Montréal will also save you some money! For \$24 per month or \$115 per season as of \href{https://bixi.com/en/pricing/}{August 2026}, you get access to electric and non-electric bikes all around the city. McGill students also get 10\% discount on seasonal memberships if you sign up early in the season (for example, \href{https://www.mcgill.ca/transport/cycling/10-discount-2026-bixi-seasonal-memberships}{2026 Bixi discount info}). Note that you need your own \textbf{helmet} to use the electric bikes. The normal ones do not require you to have any, but it's your call on safety. 

There are many \href{https://services.montreal.ca/en/maps/bike-paths}{biking streets} around Montréal, and it is highly recommended you stick to them as much as you can. Other paths will require you to be in the street and are obviously more dangerous, but are not off-limits. As a rule of thumb, what Google Maps suggests is usually a good recommendation. Also, be careful about crossing red lights since you can get a ticket for that! Same with not wearing a helmet on an electric bike. Getting a ticket this way happens more often than you might think. Lastly, headphones or headsets of any kind are prohibited on bikes.

\subsection{Car Rental}\index{Car Rental}
If you need a vehicle to move things, to explore the surrounding area, or even to go on a trip somewhere, McGill has a pretty good \href{https://www.mcgill.ca/travelservices/transport/book-vehicle/enterprise}{discount} with Enterprise. Since different companies bill you for distance, time, or fuel, it might be more cost-efficient to go with Communauto or Leo depending on the type of trip, but McGill does not have a deal with either.


\section{Exploring the City}\index{Exploring the City}\label{exploreMTL}
Prof. Katelin Schutz made a \href{https://indico.global/event/551/attachments/56549/116890/activities.pdf}{quick city guide} for a conference in 2025 you might find interesting! Sadly, it is mostly centred around the Mile-End neighbourhood because of distance and time constraints in the context of the conference.

Another great way to explore Montr\'eal and the surrounding area is by bike as the province has many scenic biking routes. Some of the best trails are: \index{Biking}
\begin{itemize}
    \item \href{https://en.ptittraindunord.com/}{P'tit Train du Nord}
    \item \href{https://www.parcjeandrapeau.com/en/circuit-gilles-villeneuve-multi-purpose-track-provelo-sports-training-montreal/}{Circuit Gilles-Villeneuve} (ie. Formula 1 track on Jean Drapeau)
    \item \href{https://gpcqm.ca/en/grand-prix-cycliste-montreal/}{Grand Prix Circuit} on Mont Royal
    \item \href{https://niche-canada.org/2014/02/12/a-river-runs-under-along-and-through-it-montreal-and-the-st-lawrence-seaway/}{St. Lawrence Seaway dyke}: A skinny bike route right in the middle of the St. Lawrence river
    \item \href{https://www.routeverte.com/carte-interactive/}{Route Verte}
        \begin{itemize}
            \item To Chambly
            \item Along the Lachine Canal to Parc Summerlea
            \item Along the water to Parc national des Îles-de-Boucherville
            \item To Parc national d'Oka
        \end{itemize}
\end{itemize}

For another set of recommendations in various categories, check out \href{https://sophiarubens.github.io/l-ile-infini/}{this} informal guide to Montr\'eal created by a McGill PhD student in August 2026. Some of the recommendations overlap with those in this guide, but there are plenty of linearly independent recommendations, particularly related to live world music, walks in outlying neighbourhoods, the metro, and independent grocery stores. There is also a Google Maps list made by another student with places and restaurants \href{https://maps.app.goo.gl/NpDyFci6MpNeW9ef7}{around the province} and \href{https://maps.app.goo.gl/VTJt8hk3R4DSWehXA}{in the city}.

\subsection{Festivals and Events}\index{Exploring the City! Festivals and Events}
If you need the general list of things to do around the city, then you should look at \href{https://www.mtl.org/en/what-to-do/festivals-and-events}{Tourisme Montréal}, it should have most, if not all, events happening around the city! 

  \noindent If you are looking for festivals, music shows and lively events, you should look into the \textbf{\href{https://www.quartierdesspectacles.com/en/festivals-and-events}{Quartier des spectacles}}. It hosts many of the city's events and usually has multiple free venues to enjoy while walking around the sector. The main parts of the \textit{Quartier des spectacles} are the \textit{Place des arts}, the \textit{Esplanade tranquille}, and the \textit{Place des festivals}, but there are \href{https://www.quartierdesspectacles.com/fr/places-publiques}{other venues}.

Here is a list of some of the most notable events:

\subsubsection{\href{https://www.mtl.org/en/experience/montreal-winter-festivals-shine-bright}{Winter}}
\begin{itemize}
    \item \textbf{Free Skating}: At Esplanade Tranquile, Lac Castor on Mont Royal, and other open rinks. Just bring your own skates or rent them for \sim ~\$10.
    \item \textbf{Nuit blanche}: A yearly event where the city comes to life for a night in the middle of winter.
    \item \textbf{Montréal en Lumière}: A double-sided festival of food and culture with art installations throughout the \textit{Quartier des spectacles}.
    \item \textbf{Igloofest}: An electronic music festival held during winter!
\end{itemize}

\subsubsection{\href{https://www.mtl.org/en/experience/spring-activities}{Spring}}
\begin{itemize}
    \item \textbf{Temps des sucres/Sugar Shacks}: Not an event \textit{per se}, but when the snow starts to melt, sugar shacks open with the new harvest of maple water and syrup, and it's a somewhat traditional thing to visit them for brunch.
    \item \textbf{MURAL Festival}: Each year, the \textit{boulevard Saint-Laurent} is painted with murals turning the street into an open-air museum.
    \item \textbf{Grand Prix}: Montréal has it's own Grand Prix, if you like F1, this could be a wonderful oportunity!
\end{itemize}

\subsubsection{\href{https://www.mtl.org/en/experience/montreal-summer-festival-guide}{Summer}}
\begin{itemize}
     \item \textbf{Osheaga}: A summer music and arts festival! Headliners are usually big names like Twenty One Pilots, Olivia Rodrigo, Lorde, etc.
    \item \textbf{Festival international de jazz de Montréal}: A Jazz festival lasting two weeks with plenty of free shows.
    \item \textbf{\textit{Les Francos}}: A Francophone music festival lasting a bit less than two weeks. Works in a similar way to the jazz festival.
    \item The \textbf{OM} (Orchestre Métropolitain) and \textbf{OSM} (Orchestre Symphonique de Montréal) do yearly free shows during the summer, they are worth looking into! The OSM also has cheap tickets ($\sim\$35$ in 2026) for people under 35 years of age, with an additional 25\% off for a set of four tickets.
    \item \textbf{Saint-Jean/Fête nationale du Québec}: The province's national holiday. There are multiple shows and cultural events in the days before the event and a single massive show on the day of.
    \item \textbf{Just For Laughs festival}: A comedy festival featuring both free and paid shows of stand-up, improv, sketch, and other types of comedy.
\end{itemize}     
\subsubsection{\href{https://www.mtl.org/en/experience/fall-festivals-guide}{Fall}}
\begin{itemize}
     \item \textbf{\textit{Jardins de lumière}}: A festival of lanterns and silk in the botanical gardens and olympic park.
    \item \textbf{Ramen Ramen Fes}: A ramen festival with restaurants celebrating ramen with their own original bowls with a ranking made by the public and judges at the end of the event.
    \item \textbf{\textit{MTLàTABLE}}: À food festival where restaurants create unique menues for about half a month.
\end{itemize}     



\subsection{Activities}\index{Exploring the City! Activities}

Here are some year-long and seasonal things to do in the city:

\subsubsection{Food}
Montréal has many food festivals and events, from the Korean and Japanese food festivals to the monthly food truck gatherings, all the way to chefs giving cooking demonstrations during the many winter events. If you are interested, feel free to look at the \href{https://www.mtl.org/en/experience/guide-montreal-food-festivals}{schedule} and plan your meals accordingly! Also, if you ever want to try new restaurants but don't have people willing to join you, we have a \href{https://chat.whatsapp.com/H7Re6U6d7sQBbr9uv3EaJo?mode=gi_t}{group chat} just for that! We go on random expeditions to try new cuisines, experience neighbourhoods, and go where we wouldn't go otherwise!
\subsubsection{Attractions}
\begin{itemize}
    \item If you are looking for \textbf{places to walk around}, here are a few places with nice shops or things to visit: Mont-Royal (both the street and the mountain), Boulevard St-Laurent, Rue St-Denis, Avenue Laurier, Rue Rachel, Parc La Fontaine (sometimes there are even \href{https://montreal.ca/evenements?dc_relation.url=%2Flieux%2Ftheatre-de-verdure&orderBy=dc_temporal.start}{events}!), Parc Jean-Drapeau, Canal Lachine, Old Port, and the whole of the Plateau and mile-end neighbourhoods.
    \item Looking for \href{https://museesmontreal.org/en/museums}{\textbf{museums}}? Visit the Biodome, Biosphere, Planetarium, Insectarium, and Botanical gardens for a science and nature-themed experience! Otherwise, the city is full of more traditional museums like the Museum of Fine Arts (that is free for people 25 and under), the Canadian Centre for Architecture, and the Musée d'art contemporain de Montréal (Montréal's museum of contemporary art).
\end{itemize}

\subsubsection{Adventure}
\begin{itemize}
    \item Just as it has grown in popularity over the last few years, there is a somewhat big \textbf{climbing/bouldering} community at McGill and in Montréal! If you wish to try it out, look for the information in the MGAPS newsletter as we often do introduction events for people new to the sport. Otherwise, Allez-Up has a membership valid for all three of their gyms, including one where you can do top-rope and lead. Other locations include Café Bloc, Bloc Shop, Nomad Bloc (outdoor bouldering in the summer), and Shakti Rock Gym. 
    \item While quite unusual, you can try \textbf{Axe Throwing} while in Montréal! There are a few places such as \textit{Rage Montréal} and \textit{Sports de combats}. They will lend you the axes and teach you how to safely practice the sport. It's a fun thing to try and great for letting out some frustration!
    \item If you want an evening out with friends without having to plan much or being constrained to a specific activity, Montréal has many board game pubs where you can enjoy a wide selection of short or long games for a reasonable fee! In no specific order, there is \href{https://maps.app.goo.gl/KFguqV2jX4RTyHjC7}{Randolph}, \href{https://maps.app.goo.gl/Tg3gvhN78JZaeNzb6}{Colonel Moutard}, and \href{https://maps.app.goo.gl/tSmz2U3ZADBaTKhc8}{Aysse}.
\end{itemize}

\subsubsection{Winter}
    \begin{itemize}
            \item \textbf{Ice Skating} is very frequent in the winter, and many places have outdoor ice-skating rinks! The McGill campus has its own, but if you don't have your own skates, the rink at \textit{Esplanade tranquille} offers rentals for a cheap price. If you wish to skate in the summer, some indoor rinks also exist. If you visit Ottawa (around 2h west of Montréal) in winter, you might also be able to skate on the Rideau Canal if weather permits (better to check \href{https://ncc-ccn.gc.ca/places/rideau-canal-skateway}{online} first).
    \item Visiting \textbf{Christmas markets} is a wonderful activity to pass the time and cheer up during the long winter nights. They have overpriced goods you will likely buy anyway and comforting food that will warm you up in the middle of the cold evening! Montréal has a few small ones, and Québec City has a big one.
    \item If you can find a few people to carshare with, \textbf{Skiing and Snowboarding} on the many mountains around the city is a great activity to spend a day away from the noise of the city and the worries of work! Depending on the time of year you read this, it might be too late for a season pass, but if you know people who have some, they might have coupons to make your day passes cheaper!
    \end{itemize}

\subsubsection{Miscellaneous}
\begin{itemize}
    \item Tuesday deals at \textbf{Cineplex} offer cheap movies every week. Assemble your group and save some money on something you've been meaning to see or enjoy a random movie you might not pay for otherwise!
    
    \item  Discounts, see \hyperref[discounts]{\textcolor{purble}{\ref*{discounts}}}.
\end{itemize}

\subsection*{Hockey culture}
Just like our distant Finnish friends, the average Canadian is most likely a (Ice) Hockey fan! It is the most popular sport, and as a new resident in the home city of the Montréal Canadiens (yes, \textit{Canadiens} is in French even when using the English name) you'll quickly get to see the people get crazy every time the \textit{Habs} win. If you wish to take a more active part in the hockey culture, some of the early-season games against... less-favoured teams might be cheap enough to justify a ticket to the Bell Centre. Otherwise, most pubs will have the game on their screens. If the Canadiens are in the playoffs, then every bar will have additional screens (sometimes even projectors) to watch the games, and they most likely will be full during games. It's a great experience to watch the team win and have everyone in the establishment jump in excitement!

The newer addition to the hockey scene is the PWHL (Professional Women's Hockey League) and its Montréal team: La Victoire. With tickets at a fraction of the price of the impressively expensive NHL games, the PWHL offers wonderful hockey plays and an incredible atmosphere for $\approx \$ 30-40$! Located in the annoyingly slightly-too-far-for-Zone-A Place Bell, the Victoire have become an integral part of the city and gather a crowd every game. Highly recommend going to watch a match with a group of friends! 



\subsection{Restaurants and Food}\index{Exploring the City! Restaurants and Food}

Without much context, as you might as well go look into the restaurants yourself and try them in person, here is a list of restaurant recommendations:

\subsubsection{Close to McGill}
\begin{table}[H]
    \centering
    \arrayrulecolor{ocre}
    \begin{tabular}{lr}
    \hline
        Vinh's (Vietnamese, only during semester) &  Boustan (Lebanese)\\
        Nouilles Zhonghua (Chinese noodles) & Opiano (Korean)\\
        Mae Sri Comptoir Thai (Thai) & Le chef (Turkish) \\
        Shawarmaz (Lebanese) & Bistro La Cite (North American Diner-style)\\
        \hline
    \end{tabular}
\end{table}

\subsubsection{Around the City}
\begin{table}[H]
    \centering
    \arrayrulecolor{ocre}
    \begin{tabular}{lr}
         \hline
         Mange dans mon hood (MDMH) (burgers and poutine) &  Le Super Qualité (Indian)\\
         Le petit coin dumpling (dumplings) & Le central (food court)\\
        Flamme du Sichuan (Sichuan cuisine) & Comon verdun (Korean) \\
        Mauve (flower shop X restaurant X wine bar) & Café Régine (breakfast)\\
        Iconoglace/Kem Coba/Ca Lem (ice cream) & Toaster Villeray (breakfast)\\
        Kan Bei (Sichuan cuisine) \\
         \hline
    \end{tabular}
\end{table}

\textit{Première Moisson} has good bread and pastries, the markets usually have good (sometimes expensive) food stalls, and local grocery stores often have better produce and sometimes lower prices.

\subsubsection{Bars}
\begin{itemize}
    \item \textbf{Bar Le Record}: A place with a cozy atmosphere with an owner that has an insane vinyl collection and plays music from it
    \item \textbf{Chez Ernest}: Absinthe bar with a very good selection
    \item \textbf{Le 4e mur}: A speakeasy style bar with a hard-to-find entrance and a cool ambiance
    \item \textbf{Bar Darling}: A classic bar, somewhat expensive, but with a nice vibe
    \item \textbf{Siboire}: More a pub than a bar, but they have nice beers on rotation!
    \item \textbf{Le terminal}: A bar with a broad selection of beers, plenty of events, and multiple stand-up nights
\end{itemize}

 \subsubsection{On the More Expensive Side}
  \begin{itemize}
    \item \textbf{Miracolo}: They have a very nice tasting menu that is very reasonable for two people!
    \item \textbf{Le Majestique}: Fancy-ish restaurant with very good food and a pleasant atmosphere
    \item \textbf{Daldongnae Korean BBQ}: A good korean BBQ, it's a fun place to go with a group of friends! 
    \item \textbf{O-Thym}: Fine dining, with the prices that go with it, so expect 100 for a 3-course meal
\end{itemize}

  \subsubsection{Farmer's Markets}
  The \textbf{\textit{Marchés publics de Montréal}} offer fresh produce and local products year-round. They are also cosy places to visit and generally have a friendly atmosphere! The Jean-Talon Market is the biggest one and is surrounded by many shops that are worth a visit! The Atwater Market is also a great place to visit, especially at the end of the calendar year when it has a Christmas market section. Many procude stalls in the markets have increadible deals on ugly/old produce, if you live nearby and are looking for things on a budget, this can certainly save you quite a lot of money!
  
  All the markets host multiple \href{https://www.marchespublics-mtl.com/en/events}{events}, from seed-planting to fermentation workshops and even traditional dance classes. If you live near one of them, you can often find cheaper and higher-quality in-season produce at these markets than at the grocery store. 

\subsubsection{Bagels}
Originating from Jewish immigration alongside Montréal smoked meat, the Montréal-style \textbf{bagels} are an undeniable part of the food scene. The two main establishments are in an everlasting popularity war, so feel free to go try them both and decide your favourite between Fairmount and St-Viateur! (Author's note: This guide is meant to be helpful and impartial, so I will admit that both are very good and I won't mention my favourite, but I want to say that there is a correct answer).

\subsection*{Poutine Tier list}

\noindent Poutine is Canada's unofficial national dish, made of crispy French fries and fresh cheese curds topped with hot brown sauce (meat stock-based sauce). As a staple from Québécois cuisine, it can easily be found in most pub-style restaurants and \textit{Casse-croûtes} (snack bars usually found on the outskirts of towns, in small villages, and on the side of the highways). 

If you visit further into the province, you might stumble upon \textit{Drummondville} and \textit{Victoriaville}, which both contest the claim for the birthplace of poutine. Up to you to decide which of these two towns makes the best one!

\noindent\textbf{\color{ocre}Disclaimer:} The views and opinions expressed here are those of the authors and do not reflect the official policy of MGAPS or McGill Physics. This list only includes poutine from Montréal and environs.

\newcommand{\tieritem}[1]{%
    \tcbox[on line, size=small, colback=white, colframe=gray!50]{#1}\;}
    
\begin{table}[hbt]
\centering
\arrayrulecolor{black}
\renewcommand{\arraystretch}{1.5}
\begin{tabular}{|>{\centering\arraybackslash}m{1cm}|m{12cm}|}
\hline
\rowcolor{red!70}\textbf{\sffamily S} & 
\tieritem{\href{https://maps.app.goo.gl/8QTYi7h8oizniiQm9}{\color{black}Alfa}} 
\tieritem{\href{https://maps.app.goo.gl/ptYJXzRPfnk5hx3f9}{\color{black}MDMH}} \tieritem{\href{https://maps.app.goo.gl/Rss6FN5J8ahoBE3t6}{\color{black}Patate Rouge}}
\tieritem{\href{https://maps.app.goo.gl/tBWkx6M7hnajEKMu6}{\color{black}Chez Claudette}}\\
\hline
\rowcolor{orange!70}  & 
\tieritem{\href{https://maps.app.goo.gl/qeHadKEwxvhLXemFA}{\color{black}Greenspot}}\tieritem{\href{https://maps.app.goo.gl/8is7FU3e59gASaFr5}{\color{black}La Pataterie} }
\tieritem{\href{https://maps.app.goo.gl/ka3uXE781Ldz7p6u7}{\color{black}Lafleur}} 
\tieritem{\href{https://maps.app.goo.gl/wZuM8Yxkfg4W3KadA}{\color{black} Ma Poule Mouillée}}\\
\rowcolor{orange!70}\multirow{-2}{*}{\textbf{\sffamily A}}&
\tieritem{\href{https://maps.app.goo.gl/h5gkmJH98gf2xxWs7}{\color{black} Souvlaki Bar}}
\tieritem{\href{https://maps.app.goo.gl/FefHQE24P1qBpAmY8}{\color{black}Le Petit Wellington}}
\tieritem{\href{https://maps.app.goo.gl/bxzypkpnvS3gv9pd8}{\color{black}Shawarmaz}}\\
\hline
\rowcolor{yellow!70}  & 
\tieritem{\href{https://maps.app.goo.gl/UnR7Y8MLVW3KZeP7A}{\color{black}Orange Julep}} 
\tieritem{\href{https://maps.app.goo.gl/jzdqsctfxLM7WuKa9}{\color{black}La Belle Province}}
\tieritem{\href{https://maps.app.goo.gl/8skjWhNfqJ57tHs29}{\color{black}Dunn's Famous}}\\ 
\rowcolor{yellow!70} \multirow{-2}{*}{\textbf{\sffamily B}}& 
\tieritem{\href{https://maps.app.goo.gl/spmNha2nXjEs23Sr8}{\color{black}La Banquise}}
\tieritem{\href{https://maps.app.goo.gl/QakCDGAxBZgNok2V7}{\color{black} 3 Brasseurs}} \tieritem{\href{https://maps.app.goo.gl/7XLJGcvbQhbJnv7e7}{\color{black} Copper Branch}} \\
\hline
\rowcolor{green!50}  \textbf{\sffamily C} & \tieritem{\href{https://maps.app.goo.gl/EGQXYtuXGiwzb4319}{\color{black}Costco}}\tieritem{\href{https://maps.app.goo.gl/uBCyVXfnzfJnt5zf6}{\color{black}Boustan}} \\
\hline
\rowcolor{gray!40}  \textbf{\sffamily F} & \tieritem{\href{https://maps.app.goo.gl/HmMBu9eKwjGAGxjx7}{\color{black}McDonald's}}\\
\hline
\end{tabular}
\end{table}

\subsection{Music Venues}
See the Montreal guide listed at the beginning of Section \hyperref[exploreMTL]{\textcolor{purble}{\ref*{exploreMTL}}}  for more info, but here's a quick overview strongly biased towards world music---resulting from the biases of the author of this section---organized into price tiers. 
\subsection*{Free-\$5}
Offerings in this price category are wonderfully plentiful. Because Montréal has such an active arts scene, if you keep your finger on the pulse, you'll often get to hear artists before they gain wider notoriety. For example, the author of this section had the chance to hear a Djely Tapa concert for free in Parc des Faubourgs during summer 2025---tickets for her concert at Place-des-Arts are now on sale!

Many festivals that set up stages in Place-des-Arts in the summer---e.g. Festival Nuits d'Afrique and JazzFest---have stages where you can listen to incredible music for free, but there is a lot more free music in this city if you know where to look. 

The indispensable guide to this tier is the \href{https://montreal.ca/calendrier-culturel?mtl_content.evenements.event_type.subType.code=TEV15&shownResults=12}{Montréal cultural calendar}, which lists many one-off and series events sponsored by all boroughs. Many host venues are public parks and Maisons de Culture. If a certain borough, concert series, cultural organization, or venue particularly resonates with you, you can usually sign up for their mailing list. 

The Syli d'Or (African and Afro-Diaspora battle of the bands which Club Balattou puts on every February-April) charges for entry to knockout rounds, but the preliminary rounds are free to attend (although coat check is \$2).

\subsection*{\$15-\$35}
If you spend long enough on the lookout for live world music in Montreal and spend any time on social media, expect to encounter targeted ads for events ticketed by Le Point de Vente. Although you might not expect a particular ticket distributor to have a seemingly curated blend of events---and there are many events on the site entirely unrelated to world music---this is an incredibly fruitful place to peruse options for chill bar concerts. 

Venues to keep an eye out for that often have interesting music offerings: Quai des Brumes, L'escogriffe, Club Soda, Bar le Ritz PDB, Club Balattou, Cabaret Lion d'Or, La Cale pub zéro dechet, Aux Angles Ronds, Casa del Popolo, Theâtre Rialto.

\subsection*{$\gtrsim$\$35}
In this price tier, you're more likely to encounter ``big name" artists and venues; Place-des-Arts, MTELUS, and the Bell Centre, to name a handful. 

\subsection*{Bonus: Cultural Organizations}
By keeping an eye out for cultural festivals, you can often discover more music, often with incredibly accessible ticket prices. To name just a few such events and organizations: Kurdish Festival, Festival Sefarad, KlezKanada Klezmer Jams, Festival Orientalys, and Casal Català de Quebec.

Blurring the lines between festivals and cultural organizations, there's another interesting category of musical offerings in this city. Two great examples are Festival Gnaoua and the offerings of le Centre des Musiciens du Monde (CMM). CMM hosts artists in residence throughout the year, but two particular highlights are Festival Racines and Musique sous un Arbre. The first includes concerts, workshops, lectures, and Q\&As with the artists, and multi-concert passes are available, meaning you can see an incredible musical diversity (for example, with a 2026 five-concert pass, the author of this section had the privilege of hearing Innu-Occitan fusion, Anatolian-Peruvian fusion, and more). The second includes free concerts in Parc Lahaie (rain location: the adjacent church) twice per week during the middle portion of the summer. 

\section{Exploring the Province}\index{Exploring the Province}
\begin{itemize}
    \item Go see nature in one of the \href{https://www.sepaq.com/carte/index.dot?language_id=1}{SÉPAQ nature reserves }
    \item Visit the Granby Zoo or the Ecomuseum Zoo.
    \item Visit the National Capital (Here is the author's personal list of recommendations for \href{https://docs.google.com/document/d/17WaLQl0T-heW3wgH_kzgjfSy8mMwNQudzKE0SNTB6VA/edit?usp=sharing}{Québec City}).
    \item Visit Baie-Saint-Paul, a charming small city a bit to the east of Québec City.
    \item Visit Tadoussac to see the whales!
    \item Visit Saguenay for beautiful nature, lakes, and camping locations.
    \item Visit Bas-Saint-Laurent (Kamouraska and Rimouski) and Gaspésie (Gaspé, Persé, and Carleton-sur-Mer are wonderful towns to visit, but the whole coastline is beautiful!). If you pass by Gaspésie and like smoked fish, you might want to drop by \textit{Atkins et Frères} in Mont-Louis.
    \item This is the same list as above, but here is a Google Maps list of things \href{https://maps.app.goo.gl/NpDyFci6MpNeW9ef7}{around the province}.
\end{itemize}